\subsection{Постановка задачи и результат работы}
Требуется классифицировать изображения пяти видов цветов: ромашки, одуванчики, розы, подсолнухи и тюльпаны. Исходное задание использует набор \href{https://www.kaggle.com/datasets/alxmamaev/flowers-recognition}{Flowers Recognition}. Изображения помещаются в каталоги \file{daisy}, \file{dandelion}, \file{rose}, \file{sunflower}, \file{tulip}; число файлов и соответствие индексов меткам нужно получить из своей копии данных.

\textbf{Обязательное сравнение сохраняет три части исходника:}
\begin{enumerate}
\item заданная свёрточная сеть, обучаемая с нуля;
\item предобученная свёрточная архитектура с новой головой классификации;
\item собственная архитектура студента, обучаемая с нуля.
\end{enumerate}
Необходимо объяснить устройство моделей, визуализировать результаты и сравнить качество при одинаковом разбиении. В пособии используется PyTorch, как и в коде исходного ноутбука. При выполнении на TensorFlow/Keras нужно явно перенести слои, режимы, предобработку и протокол, а не смешивать API двух фреймворков.

\textbf{Результат сдачи:} схема каждой архитектуры, число всех и обучаемых параметров, кривые обучения, таблица сравнения трёх моделей, матрицы ошибок, примеры прогнозов, сохранённые веса и описание воспроизводимого запуска.

\subsection{Изображение как тензор}
Тензор здесь --- многомерный массив чисел. Цветное изображение содержит три канала RGB: красный, зелёный и синий. В PyTorch пакет изображений имеет форму \((B,C,H,W)\): размер пакета, число каналов, высота, ширина. Для принятого входа это \((B,3,224,224)\). \file{ImageFolder} читает метки из имён каталогов и формирует \file{class\_to\_idx}; используйте это отображение и при обучении, и при выводе названий~\cite{imagefolder}.

Преобразование в тензор переводит обычные 8-битные RGB-пиксели в значения от 0 до 1. Затем выполняется нормализация каналов:
\begin{equation}
z_{c,h,w}=\frac{x_{c,h,w}-\mu_c}{\sigma_c},\qquad
x_{c,h,w}=\sigma_c z_{c,h,w}+\mu_c.
\end{equation}
В исходном примере используются \(\mu=(0{,}485,0{,}456,0{,}406)\) и \(\sigma=(0{,}229,0{,}224,0{,}225)\). Это фиксированные константы нормализации ImageNet, а не статистики, вычисленные по цветам. Для визуализации сначала выполните обратное преобразование; прямой \file{clip(0,1)} нормализованного тензора искажает цвета.

В исходном сравнении все модели получают \file{Resize((224,224))}, преобразование в тензор и одну нормализацию. Такое изменение размера может менять пропорции изображения. У весов ResNet18 есть собственный рекомендованный рецепт \file{weights.transforms()}: изменение размера и центральная обрезка~\cite{resnet18}. Это другой рецепт. Если сравниваете его с исходным, оформите отдельный опыт и применяйте выбранные правила одинаково к сравниваемым моделям.

\textbf{Аугментация} создаёт изменённые варианты обучающих изображений, например зеркальные отражения. Она помогает показать модели больше допустимых вариантов объекта, но преобразование должно сохранять его класс. В основном сравнении этой работы случайные преобразования выключены; их влияние можно исследовать отдельно. На валидации применяются фиксированные преобразования.

\subsection{Свёрточный слой}
Полносвязный слой рассматривает все входные координаты совместно. Свёрточный слой использует небольшой локальный участок и одни и те же веса в разных пространственных положениях. Для шага 1 и без дополнения границы один выходной канал вычисляется по правилу
\begin{equation}
a_{o,h,w}=b_o+\sum_{c=1}^{C_{in}}\sum_{u=0}^{k_h-1}\sum_{v=0}^{k_w-1}
W_{o,c,u,v}\,x_{c,h+u,w+v}.
\end{equation}
В \file{Conv2d} реализуется взаимная корреляция: ядро в этой записи не переворачивается. В терминологии нейронных сетей операцию традиционно называют свёрткой~\cite{conv2d}. Полученные карты признаков --- таблицы откликов фильтров в разных участках изображения --- проходят через нелинейность, например ReLU. Обучение определяет, какие локальные сочетания полезны для классификации.

\begin{keybox}{Маленький расчёт}
Для одной карты \(X=\left[\begin{smallmatrix}1&2&3\\4&5&6\\7&8&9\end{smallmatrix}\right]\), ядра \(K=\left[\begin{smallmatrix}1&0\\0&2\end{smallmatrix}\right]\), нулевого смещения и шага 1 без padding получается \(\left[\begin{smallmatrix}11&14\\20&23\end{smallmatrix}\right]\). Например, левый верхний элемент равен \(1+2\cdot5=11\). Найдите остальные элементы вручную.
\end{keybox}

Для размера входа \(H\), ядра \(k\), шага \(s\), padding \(p\) и dilation \(d\) высота выхода равна
\begin{equation}
H_{out}=\left\lfloor\frac{H+2p-d(k-1)-1}{s}+1\right\rfloor.
\end{equation}
Для ширины действует такая же формула. В этой работе \(d=1\). При \(k=3\), \(s=1\), \(p=1\) размер сохраняется. Для обычной свёртки без групп со смещением число параметров
\begin{equation}
P=C_{out}(C_{in}k_hk_w+1).
\end{equation}
Первый слой \(3\to32\) с ядром \(3\times3\) содержит \(32(3\cdot3\cdot3+1)=896\) параметров. Общие веса резко уменьшают число параметров относительно полносвязного соединения всех пикселей со всеми пространственными выходами.

\subsection{Pooling и устройство исходной сети}
\file{MaxPool2d(2,2)} выбирает максимум в каждом окне \(2\times2\) и уменьшает высоту и ширину вдвое для чётного входа. Он не обучает веса. Pooling уменьшает пространственную детализацию и стоимость следующих вычислений; это не означает полной устойчивости классификатора к любому сдвигу или повороту.

\begin{table}[H]\centering\small
\caption{Размерности исходной CNN при входе 224\(\times\)224}
\begin{tabularx}{\textwidth}{@{}p{0.58\textwidth}Y@{}}
\toprule
Последовательность операций & Форма без размера пакета \\
\midrule
Вход RGB & \(3\times224\times224\) \\
Conv \(3\to32\), ReLU, Conv \(32\to64\), ReLU, MaxPool & \(64\times112\times112\) \\
Conv \(64\to128\), ReLU, Conv \(128\to128\), ReLU, MaxPool & \(128\times56\times56\) \\
Conv \(128\to256\), ReLU, Conv \(256\to256\), ReLU, MaxPool & \(256\times28\times28\) \\
Flatten & \(200704\) \\
Linear, ReLU; Linear, ReLU; Linear & \(1024\to512\to5\) \\
\bottomrule
\end{tabularx}\end{table}

Все свёртки в таблице имеют ядро 3, шаг 1 и padding 1. Полная реализация \file{original\_cnn()} в учебном модуле сохраняет размеры слоёв полученного ноутбука. Flatten только меняет представление тензора: он не извлекает новые признаки.

\begin{keybox}{Почему нужно заранее оценить память}
Слой \(200704\to1024\) содержит \(205\,521\,920\) параметров со смещением. Вся исходная сеть --- \(207\,175\,365\). Только веса float32 занимают около 0,77 GiB. Веса, градиенты и два состояния Adam суммарно требуют около 3,09 GiB, ещё без активаций и временных буферов. Это расчёт по архитектуре, а не измерение пикового расхода памяти.
\end{keybox}

Для собственной модели полезно исследовать глобальное усреднение карт. В примере оно выполняется слоем \file{AdaptiveAvgPool2d(1)}. Этот слой превращает \(C\times H\times W\) в \(C\times1\times1\), после чего компактная голова предсказывает пять классов. Это уменьшает число параметров, но меняет модель и не гарантирует лучшего качества. В комплекте дан небольшой пример; самостоятельная архитектура должна содержать обоснованное собственное изменение.

\subsection{Логиты, cross-entropy и обучение}
Последний слой выдаёт пять логитов \(z_1,\ldots,z_5\) --- оценок классов, которые могут быть любыми вещественными числами. Softmax превращает их в положительные числа с суммой 1. Вероятности softmax определяются как
\begin{equation}
p_c=\frac{e^{z_c}}{\sum_{j=1}^{5}e^{z_j}},\qquad
L=-\frac1B\sum_{i=1}^{B}\log p_{i,y_i}.
\end{equation}
\file{CrossEntropyLoss} принимает исходные логиты формы \((B,5)\) и целочисленные индексы классов формы \((B,)\) с типом \file{torch.long}. Отдельный Softmax перед этой функцией не нужен~\cite{crossentropy}. Для выбора класса достаточно \file{logits.argmax(dim=1)}; наибольший логит и наибольшая вероятность имеют один индекс. Вероятность softmax сама по себе не является проверенной уверенностью модели.

Итерация обучения включает обнуление градиентов, прямой проход, вычисление потери, обратный проход и шаг оптимизатора. При проверке обратный проход и шаг отсутствуют. \file{model.eval()} переключает поведение слоёв, а отключение градиентов снижает затраты на вычислительный граф: это разные операции.

Метрики считают заново для каждой эпохи и каждой фазы. Средняя потеря должна учитывать размер последнего пакета:
\begin{equation}
L_{epoch}=\frac{\sum_b n_b L_b}{\sum_b n_b}.
\end{equation}
Например, для пакетов из 32 и 8 изображений нельзя брать простое среднее двух потерь. Accuracy считают по всем правильным ответам фазы, macro-F1 --- по всем её меткам и прогнозам, а не как невзвешенное среднее F1 пакетов.

\subsection{Перенос обучения и заморозка}
Предобученная сеть уже получила веса на другой задаче. Основную часть сети, преобразующую изображение в признаки, называют \textbf{извлекателем признаков (backbone)}. Завершающий блок, переводящий эти признаки в ответы нашей задачи, называют \textbf{головой классификации (head)}. В переносе обучения её признаки применяются к новому набору. В исходном варианте оставляют фиксированный извлекатель признаков и обучают новую голову; при \textbf{дообучении (fine-tuning)} дополнительно изменяют часть предобученных слоёв. Эти режимы нужно различать в отчёте~\cite{transfer}.

Пример использует ResNet18 с явно указанными весами \file{IMAGENET1K\_V1}. Основой сети служат остаточные блоки: к преобразованию входа добавляется путь, передающий сам вход. Если размеры совпадают, блок имеет вид \(h_{next}=\operatorname{ReLU}(F(h)+h)\); при изменении размеров применяется согласующее преобразование короткого пути. Четыре группы блоков используют 64, 128, 256 и 512 каналов; после глобального усреднения заменяется слой \file{fc} на \(512\to5\)~\cite{resnetpaper,resnet18}.

\file{requires\_grad=False} исключает параметры из расчёта их градиентов. Новая голова создаётся после заморозки и остаётся обучаемой; её параметры включаются в оптимизатор. BatchNorm нормализует промежуточные значения в сети. У него есть также текущие средние и дисперсии --- сохраняемые служебные значения (буферы), изменяющиеся в режиме обучения~\cite{batchnorm}. Для полностью фиксированного извлекателя в данном протоколе backbone работает в \file{eval}, а новая голова --- в \file{train}. Простого общего вызова \file{model.train()} здесь недостаточно.

Сравнение с обучением с нуля учитывает дополнительный ресурс: веса ResNet получены на внешних данных. Поэтому вывод относится к практическим стратегиям обучения при указанном предобучении; это не сравнение одной лишь архитектуры при одинаковом объёме исходной информации.
